\section{Analysis of the Problem Context}
\label{sec:updates}
% Updating a stream database involves the following two operations: deleting a stream and inserting a stream. The following section describes in detail for each operation which
% information needs to be updated in the new set of descriptive queries.

In this section, we study the problem of event
query discovery under updates, focusing especially on
the relation of the sets of matching queries before
and after an update. We first characterize the problem
context. Let $D_0=\{S_1,\ldots, S_d\}$ be a stream
database,
$\theta_0$ be a support threshold, and
$Q_{D_0}^{\theta_0}$ be a matching query set. For this
setting, we distinguish two possible updates, as
follows, which may also be combined:
\begin{itemize}[nosep,left=.2em .. 1.5em]
  \item The stream database $D_0$ is updated resulting
  in a new database $D_1$. Here, we consider the
  addition of $k$ new
  streams, i.e., $D_1 = D_0 \cup \{S_{d+1}, \dots,
  S_{d+k}\}$, as well as the removal of $k$ existing
  streams, i.e., $D_1 = D_0 \setminus D^-$ with $D^-
  \subseteq D_0$ and $|D^-|=k$.
  \item The support threshold $\theta_0$ is updated
  resulting in a new threshold $\theta_1 \in (0,1]$.
%  \item Both, the stream database $D_0$ and the
%  support threshold $\theta_0$ are updated.
\end{itemize}
These updates potentially have implications for the
set of matching queries, $Q_{D_0}^{\theta_0}$, which,
hence, needs to be updated to $Q_{D_0}^{\theta_1}$,
$Q_{D_1}^{\theta_0}$, or $Q_{D_1}^{\theta_1}$,
respectively. Specifically, we seek to
answer the following questions:
\begin{enumerate}[nosep,left=.2em .. 1.5em]
  \item Which queries are still matching? That is,
  which queries $q\in Q_{D_0}^{\theta_0}$ do still belong to
  $Q_{D_0}^{\theta_1}$,
  $Q_{D_1}^{\theta_0}$, or $Q_{D_1}^{\theta_1}$,
  respectively?
  \item Which new queries are matching? That is,
    which queries $q\notin Q_{D_0}^{\theta_0}$ belong to
    $Q_{D_0}^{\theta_1}$,
    $Q_{D_1}^{\theta_0}$, or $Q_{D_1}^{\theta_1}$,
    respectively?
\end{enumerate}


%\begin{problem}[Event Query Discovery with Updates]
%  \label{problem-updates}
%    Given a stream database $D_0=\{S_1,\ldots,
%S_n\}$, a support threshold $\theta_1$,
%    and a query set $Q_{D_0}^{\theta_0}$, and an
%updated discovery problem.
%    The possible updates include:
%    \begin{itemize}[nosep,left=1em]
%      \item The stream database $D_0$ is updated with
%$k$ new
%      streams $\{S_{d+1}, \dots, S_{d+k}\}$
%      or $k$ existing streams are removed.
%      \item The support threshold $\theta_0$ is
%updated.
%      \item Both the stream database $D_0$ and the
%support threshold $\theta_0$ are updated.
%    \end{itemize}
%    The goal is to update the query set
%$Q_{D_0}^{\theta_0}$ to reflect the changes
%    in the stream database $D_0$ and the support
%threshold $\theta_0$.
%    \begin{enumerate}[nosep,left=1em]
%      \item Which queries in $Q_{D_0}^{\theta_0}$ are
%still correct?
%      \item Which new queries are correct?
%      \item Which information can be derived from the
%previous query discovery?
%      \item How can the query set
%$Q_{D_1}^{\theta_1}$
%      be updated efficiently?
%      \item Is it always feasible to update the query
%      set $Q_{D_1}^{\theta_1}$ incrementally
%      or are we faster by simply running the query
%      discovery on the new setting?
%    \end{enumerate}
%  \end{problem}

%In this section we want to give an answer to the
%question 1-3. Question 4 will
%be answered in \autoref{sec:algos} where we present
%our algorithms.
%Finally, in \autoref{sec:evaluation} we give an
%answer to question 5 comparing
%the runtimes from a baseline event query discovery
%implementation to our
%algorithms that take into account previous
%information.

%\subsection{Relations between Matching Query Sets}
Ideally, the above questions shall be answered
incrementally. That is, we would like to avoid
computing $Q_{D_0}^{\theta_1}$,
$Q_{D_1}^{\theta_0}$, or $Q_{D_1}^{\theta_1}$ from
scratch, but rather derive it directly from
$Q_{D_0}^{\theta_0}$.

Before turning to algorithmic solutions to the above
questions, we study the three relations between the
matching query sets as
illustrated in \autoref{fig:matching_sets}.

\begin{figure}[h]
  %\begin{center}
\centering
\begin{tikzpicture}[scale=2]
    \def\is{1.0} % inner seperation of nodes

    % GROUPS
    \node[inner sep=\is] (Q00) at (0,0) {$Q_{D_0}^{\theta_0}$};
    \node[inner sep=\is] (Q01) at (0,1) {$Q_{D_0}^{\theta_1}$}; % SO+(1,3)
    \node[inner sep=\is] (Q10) at (1,0) {$Q_{D_1}^{\theta_0}$};
    \node[inner sep=\is] (Q11) at (1,1) {$Q_{D_1}^{\theta_1}$};


    % ARROWS
    \draw[<->,thick,myblue] (Q00) -- (Q01)
      node[midway,right=-.65,scale=1.1] {(a)}
      ;
    \draw[<->,thick,myyellow,dotted] (Q00) -- (Q10)
      node[midway,above=-.5,scale=1.1] {(b)}
      ;
    \draw[<->,thick,myblue] (Q10) -- (Q11)
      node[midway,left=-.65,scale=1.1] {(a)}
      ;
    \draw[<->,thick,myyellow, dotted] (Q01) -- (Q11)
      node[midway,below=-.5,scale=1.1] {(b)}
      ;
    \draw[<->,thick,mymagenta, dashdotted] (Q00) -- (Q11)
      node[pos=.75,above =-.75,scale=1.1, rotate=45]
        {(c)}
      ;
  \end{tikzpicture}
  \caption{Illustration of matching query sets and the
  relations between them.}
  \label{fig:matching_sets}
%\end{center}
\end{figure}

In this paper, we focus on the case that $D_1$ is derived from
$D_0$
by the addition of $k$ streams, i.e., $D_1=D_0 \cup
\{S_{d+1}, \dots, S_{d+k}\}$, and that the support
threshold is strengthened, i.e., $0 < \theta_0 \leq
\theta_1 \leq 1$. The reason being that the properties
discussed below hold analogously for the mirrored
updates, i.e., the deletion of existing streams and a
relaxed support threshold.
In the remainder, we first study the relation of the
sets of matching queries
when only the support threshold is updated (a), and
when only the stream database is updated (b).
Afterwards, we turn to the case that both are updated
at the same time (c).


%\textbf{Remark.} Here we only proof the cases that
%streams are added to the
%stream database and the support threshold is
%increased. But the
%porperties hold analogously for the reverted
%scenarios, that is deleting
%streams from the stream database and decreasing the
%support threshold.

%Let $D_0=[S_1,\dots,S_n]$ be a stream database and
%let $D_1$ be an updated stream database
%with $D_1=D_0 \cup \{S_{d+1}, \dots, S_{d+k}\}$.
%    Let the support thresholds $\theta_0, \theta_1$
%be fixed and $0 < \theta_0 \leq \theta_1 \leq 1$ and
%    let $s(|D_0|,\theta_0),s(|D_1|,\theta_1)$ be the
%stream thresholds for $D_0$ and $D_1$, respectively.
%    Let $Q_{D_0}^{\theta_0}$ be the set of matching
%queries for $D_0$ and $\theta_0$ and
%$Q_{D_1}^{\theta_1}$
%    be the set of matching queries for $D_1$ with
%$\theta_1$.

%In \autoref{fig:matching_sets} we have illustrated
%the matching query sets for the
%old and the updated discovery problem.

\subsection*{(a): Updating only the Support Threshold}

In case only the support threshold is changed, i.e.,
it is strengthened, we note that there exists a
containment relation between the sets of matching
queries. Based on \autoref{prop:subset}, we conclude
that the set
of matching queries for the increased support
threshold is a
subset of the one for the lower support threshold, i.e.,
$Q_{D_i}^{\theta_1} \subseteq Q_{D_i}^{\theta_0}$, $i\in \{0,1\}$.
%and $Q_{D_1}^{\theta_1} \subseteq Q_{D_1}^{\theta_0}$.

\subsection*{(b): Updating only the Stream Database}
\label{subsubsec:stream_update}
% \todo[inline]{RS: Add proof for lower bound to better explain the visualization at the end of Section 3b.}
Regarding updates of the stream database, i.e., the
addition of streams, we make the
following observation based on the stream
thresholds, i.e., the minimal numbers of streams that
need to match to fulfill the support threshold.
Consider the case that $s(|D_1|,\theta_i) =
s(|D_0|,\theta_i) + k$, $i\in \{0,1\}$, where
$k$ represents the number of added streams.
Intuitively, this means that
\emph{all} newly added streams need to be matched for
a query to remain in the set of matching queries,
under a worst-case assumption. That is, all queries
that \emph{only} match the minimal numbers of streams
to fulfill the support threshold need to match all the
new streams. We note though that, with a support
threshold of $\theta_i=1$, this worst case applies to
all queries in the set of matching queries.
If $s(|D_1|,\theta_i) =
s(|D_0|,\theta_i) + k$, then the set of
matching queries for
the updated stream database can only become smaller,
i.e., it is
a subset of the original set. This observation on a
containment relation between the matching sets of
queries under updates to the stream database is
captured by the following property:
%$Q_{D_1}^{\theta_0}
%\subseteq Q_{D_0}^{\theta_0}$.
%Analogously, we consider the case that
%$s(n+1,\theta_i)
%= s(n,\theta_i)$, $i\in \{0,1\}$, i.e., not even \emph{one} of the newly
%added streams need to be matched for a query to remain in the set of
%matching queries, the
%respective containment relation holds true:
%$Q_{D_1}^{\theta_1} \subseteq Q_{D_0}^{\theta_1}$.
%This can be summarized in the property:
\begin{property}
\label{prop:db_change_inclusion}
$s(|D_1|,\theta_i) = s(|D_0|,\theta_i) + k, i\in\{0,1\} \implies $
$Q_{D_1}^{\theta_i} \subseteq Q_{D_0}^{\theta_i}$.
\end{property}
The proof of this property is directly covered by \autoref{lemma:update},
which captures a more general case as detailed below. For now, we illustrate
what happens if the respective relation between the stream thresholds does
not hold true with an example.

Consider the streams $S_1, S_2, S_3$, as given in
\autoref{tab:example_stream_database}.
Let $D_0=\{S_1,S_2\}$ be the initial stream database and let the support
threshold be $\theta_0 = 0.6$. Then, the stream threshold is
$s(D_0,\theta_0) = 2$, i.e., a query needs to match both streams to be
frequent. Now, let $D_1=\{S_1,S_2,S_3\}$ be the stream database after adding
one stream, $S_3$, while leaving the support threshold unchanged. Then, the
stream threshold is still $s(D_1,\theta_0) = 2$ and, hence,
$s(D_1,\theta_0)\neq s(D_0,\theta_0) +k$ as $k=1$ in our example.

As mentioned above, the query $q = \langle
(x_0,\text{schedule},\text{low}),(\_,\text{schedule},\_),(x_0,\text{evict},\text{low})
\rangle$ matches $S_2$ and $S_3$, but not $S_1$.
This means that $q \in Q_{D_1}^{\theta_1}$ and $q \notin
Q_{D_0}^{\theta_1}$. As such, query $q$ illustrates that the containment
relation between the sets of matching queries does not hold and we conclude
that, if $s(D_1,\theta_0)\neq s(D_0,\theta_0) +k$, then the relation between
$Q_{D_1}^{\theta_1}$ and $Q_{D_0}^{\theta_1}$ depends on the added stream(s).


%What happens if $s_1\neq s_0 +k$?
%
%  Let $D_1=[S_1,S_2]$ be the stream database with the support threshold
%  $\theta_0 = 0.6$ and the stream threshold $s_0 = 2$.
%Furthermore, let $D_2=[S_1,S_2,S_3]$ be the neighboring stream database
%with
%the same support threshold and
%the stream threshold $s_1 = 2$.


%The following query $q = \langle
%(x_0,\text{schedule},\text{low}),(\_,\text{schedule},\_),(x_0,\text{evict},\text{low})
% \rangle$ matches $S_2$ and $S_3$, but not $S_1$.
%This means that $q \in P_{D_2}^{\theta_1}$ and $q \notin
%P_{D_1}^{\theta_1}$.
%
%Conclusion: If $s_1\neq s_0+k$, the relation of the two matching query sets
%depends on the added streams.

%\exampleboxh{Example 1}{
%What happens if $s_1\neq s_0 +k$?\\
%\label{ex1}
%Consider the following three streams
%$S_1, S_2, S_3$:
%
%	\begin{center}
%	\medskip
%	{\footnotesize
%		\begin{tabular}{l lrrr}
%			\toprule
%			Stream & Event &  Job Id  &  Status  &  Priority\\
%			\midrule
%			$S_1$ & $e_{11}$   & \textbf{\color{cyan}1} &
%\textbf{\color{cyan}schedule} & \textbf{\color{cyan}low}  \\
%			 & $e_{12}$   & 1 & kill & low  \\
%			 & $e_{13}$   & 1 & \textbf{\color{cyan}schedule} & high  \\
%			 & $e_{14}$   & \textbf{\color{cyan}1} &
%\textbf{\color{red}finish} & \textbf{\color{red}high}  \\
%			\midrule
%			$S_2$ & $e_{21}$   & \textbf{\color{cyan}2} &
%\textbf{\color{cyan}schedule} & \textbf{\color{cyan}low}  \\
%			 & $e_{22}$   & 3 & \textbf{\color{cyan}schedule} & high  \\
%			 & $e_{23}$   & \textbf{\color{cyan}2} &
%\textbf{\color{cyan}evict} & \textbf{\color{cyan}low} \\
%			\midrule
%			$S_3$ & $e_{31}$   & 4 & schedule & high  \\
%			 & $e_{32}$   & \textbf{\color{cyan}5} &
%\textbf{\color{cyan}schedule} & \textbf{\color{cyan}low}  \\
%			 & $e_{33}$   & 4 & finish & high  \\
%			 & $e_{34}$   & 6 & \textbf{\color{cyan}schedule} & low  \\
%			 & $e_{35}$   & \textbf{\color{cyan}5} &
%\textbf{\color{cyan}evict} & \textbf{\color{cyan}low}\\
%			\bottomrule
%	\end{tabular}}
%	\end{center}
%
%	\medskip
%  Let $D_1=[S_1,S_2]$ be the stream database with the support threshold
%$\theta_0 = 0.6$ and the stream threshold $s_0 = 2$.
%  Furthermore, let $D_2=[S_1,S_2,S_3]$ be the neighboring stream database
%with the same support threshold and
%  the stream threshold $s_1 = 2$.
%
%
% The following query $q = \langle
%
%(x_0,\text{schedule},\text{low}),(\_,\text{schedule},\_),(x_0,\text{evict},\text{low})
% \rangle$ matches $S_2$ and $S_3$, but not $S_1$.
% This means that $q \in P_{D_2}^{\theta_1}$ and $q \notin
%P_{D_1}^{\theta_1}$.
%
%Conclusion: If $s_1\neq s_0+k$, the relation of the two matching query sets
%depends on the added streams.
%}


The above property and our example illustrate that there is a close relation
between the amount of change applied to a stream database, i.e., the number
$k$ of added streams, the stream threshold, and the support threshold. In
fact, the support threshold can be bounded based on the other two
parameters. These bounds, in turn, enable us to derive insights on how
often we can expect \autoref{prop:db_change_inclusion} to be applicable and,
thereby, conclude that the set of matching queries of the updated database
is a subset of the original set of matching queries.

As above, let $D_0=\{S_1,\dots,S_d\}$ and $D_1=D_0 \cup \{S_{d+1}, \dots,
S_{d+k}\}$ be the original and updated databases and let $\theta_0$ be the
support
threshold for both of them. To simplify the presentation, we use the
short-hands $s_0 = s(|D_0|,\theta_0)$ and $s_1=s(|D_1|,\theta_0)$ for the
stream thresholds $D_0$ and $D_1$, respectively.
Then, it holds that:
$$
s_1 = \lceil (d+k) \cdot \theta_0 \rceil \leq \lceil d \cdot \theta_0 \rceil
+ \lceil k \cdot \theta_0 \rceil \leq \lceil d \cdot \theta_0 \rceil + k =
s_0 + k.
$$
Based thereon, and considering the fact
that $s_0 \leq s_1 \leq s_0+k$, we conclude that $s_1 \in \{s_0,\dots, s_0
+k\}$.

Moreover, knowing that $s_0-1 < \theta_0 \cdot d \leq s_0$, we derive the
following bounds for the support threshold based on the stream threshold
$s_0$ of $D_0$:
$$
\frac{s_0-1}{d} < \theta_0 \leq \frac{s_0}{d}.
$$
Analogously, we derive bounds for the support threshold $\theta_0$ based on
the stream threshold $s_1$ of the updated database $D_1$:
$$
\frac{s_1-1}{d+k} < \theta_0 \leq \frac{s_1}{d+k}.
$$
Combining the above statements, we derive bounds for the support threshold
based on the stream threshold of the original database. However,
we need to distinguish two cases, i.e., $s_1 = s_0+k$, which means that all
added streams need to be matched for a query to remain
frequent under the worst-case assumption that it only
matches exactly $s_0$ streams; and $s_1 <
s_0 +k$, which means that only some of the new streams
need to be matched for any query to remain frequent:
\begin{property}
\label{prop:db_change_bounds}
$s_1 < s_0 +k \implies \frac{s_0-1}{d} < \theta_0 \leq \frac{s_0+k-1}{d+k}$
and $s_1 = s_0+ k \implies \frac{s_0+k-1}{d+k} < \theta_0 \leq
\frac{s_0}{d}.$
\end{property}

% \begin{table}[h!]
%   \centering
%   \begin{tabular}{@{}lllllllll@{}}
%   \toprule
%            & \multicolumn{2}{l}{$s_1$:} & \multicolumn{2}{l}{$s_2$:} & \multicolumn{2}{l}{$s_2=s_1$} & \multicolumn{2}{l}{$s_2=s_1+1$} \\ \midrule
%   $s$/$s_2$ & low         & up         & low         & up          & low          & up          & low          & up            \\
%   1        & 0           & 0.25       & 0           & 0.2         & 0            & 0.2         & 0.2          & 0.25          \\
%   2        & 0.25        & 0.5        & 0.2         & 0.4         & 0.25         & 0.4         & 0.4          & 0.5           \\
%   3        & 0.5         & 0.75       & 0.4         & 0.6         & 0.5          & 0.6         & 0.6          & 0.75          \\
%   4        & 0.75        & 1          & 0.6         & 0.8         & 0.75         & 0.8         & 0.8          & 1             \\
%            &             &            & 0.8         & 1           &              &             &              &               \\ \bottomrule
%   \end{tabular}
%   \caption{The lower and upper bounds for the support threshold $\theta_1$ in relation to the stream threshold $s_1$ and $s_2$ for $n=4$.}
%   \label{tab:bounds}
%   \end{table}

% In \autoref{tab:bounds} we have illustrated the bounds for the support threshold $\theta_1$ in relation to
% the stream threshold $s_1$ and $s_2$ for $n=4$.

% The probability that for a given stream database $D_1$ with $|D_1| = n$ and a given
% stream threshold $s_1$ it holds that $s_2 = s_1$ can be calculated as follows:
% $$
% \mathbb{P}[s_2 = s_1 |n,s_1] = \frac{s_1}{n+1} - \frac{s_1-1}{n} = \frac{1}{n} - \frac{s_1}{n(n+1)}.
% $$
% And over all possible stream thresholds $s_1$ we obtain:
% % $$
% % P(s_2 = s |n) = \sum_{s=1}^{n} P(s_2 = s |n,s) = \sum_{s=1}^{n} \left( \frac{1}{n} - \frac{s}{n(n+1)} \right) =
% % 1 - \frac{1}{n(n+1)} \sum_{s=1}^{n} s = 1 - \frac{1}{n(n+1)} \frac{n(n+1)}{2} = \frac{1}{2}.
% % $$

% \begin{align*}
%     \mathbb{P}[s_2 = s_1 |n] = \sum_{s=1}^{n} \mathbb{P}[s_2 = s_1 |n,s_1] &= \sum_{s_1=1}^{n} \left( \frac{1}{n} - \frac{s_1}{n(n+1)} \right) \\
%     &= 1 - \frac{1}{n(n+1)} \sum_{s_1=1}^{n} s_1 \\
%     &= 1 - \frac{1}{n(n+1)} \frac{n(n+1)}{2} = \frac{1}{2}.
% \end{align*}

% In other words, the probability that the stream threshold $s_2$ is equal to the stream threshold $s_1$ is $\frac{1}{2}$.
% This also means that the probability that the stream threshold $s_2$ is equal to the stream threshold $s_1+1$ is $\frac{1}{2}$.

The bounds as stated in \autoref{prop:db_change_bounds} are useful to derive
insights in the applicability of
\autoref{prop:db_change_inclusion}, i.e., they enable
us to understand when $s_1=s_0+k$ holds after $k$ streams
have been added to a stream database. Intuitively, the
bounds on $\theta_0$ characterize the change points
for the stream threshold, depending on the size of a
stream database $d$ and the size of the change $k$.
We visualize the change points as determined by
\autoref{prop:db_change_bounds} in
\autoref{fig:thresholds}. Here, we consider databases
of different sizes, $d\in \{10,50,100,500,1000\}$, to
which changes of size $k\in \{1,2,5\}$ are applied.
The orange ranges characterize the intervals of the
support threshold $\theta_0$ for which it holds that
$s_1 = s_0+k$, i.e., for which we know that the set of
matching queries after the change is a subset of the
original set of matching queries.
The lower bound of the orange range is determined by
the size of the database and the size of the change,
and can be derived from \autoref{prop:db_change_bounds} as follows:
$$
\frac{s_0}{d}> \frac{s_0+k-1}{d+k} \implies \frac{s_0}{d}>\frac{k-1}{k} \implies \theta_0>\frac{k-1}{k}.
$$
From the visualization, we draw two main conclusions:
\begin{itemize}[nosep,left=.2em .. 1.5em]
  \item The higher the support threshold, the more
  likely
  it is that \autoref{prop:db_change_inclusion} is
  applicable, so that the containment relation between
  the sets of matching queries holds true.
  \item The
  larger the size of the change, the less likely the
  property is applicable.
\end{itemize}





%based on the size of the database, make an estimation
%änderung des stream thresholds in abhängigkeit von
%datenbankgröße und support value
%applicability: wann ist
%abhängig vom support und vom verhältnis der größe
%Bounds auf theta erlauben Rückschlüsse für welche
%theta man eine änderung in dem stream threshold hat
%gegeben datenbankgröße, gegeben ein k
%Daraus folgt die Anwendbarkeit von Property 2
\begin{figure}[h]
  \centering
  \includegraphics[width=.75\textwidth]{img/support_distribution_all.pdf}
  \caption{Visualization of the change points for the
  stream threshold in relation to the support
  threshold, for different scenarios of database sizes
  and change sizes.}
  \label{fig:thresholds}
\end{figure}



% But the probability distribution of the support threshold $\theta_1$ in relation to the stream threshold $s_1$ is not uniform.


% There is a clear trend that the higher the support threshold $\theta_1$
% the more likely it is to fulfill the condition $s_2 = s_1+1$.
% Let_2s consider the case $s_2 = s$.
% Then it holds:
% $$
% \lceil (n+1) \cdot \theta_1 \rceil < \lceil n \cdot \theta_1 \rceil + 1.
% $$
% From this we can deduce that $n \cdot \theta_1 + \theta_1 \leq s$.
% Therefore, $\theta_1 \leq \frac{s}{n+1}$.
% Analogously, we obtain for the case $s_2 = s+1$: $\theta_1 > \frac{s}{n+1}$.



% Then
% \begin{align*}
%     \theta_1 &\leq \frac{\lceil n \cdot \theta_1 \rceil}{n+1} \\
%     (n+1) \cdot \theta_1 &\leq \lceil n \cdot \theta_1 \rceil\\
%     \theta_1 \cdot n + \theta_1 &\leq \lceil n \cdot \theta_1 \rceil \\
%     \theta_1 &\leq \lceil n \cdot \theta_1 \rceil - n \cdot \theta_1.
% \end{align*}
% $$
% \theta_1 \leq \frac{\lceil n \cdot \theta_1 \rceil}{n+1} = \frac{\lceil n \cdot \theta_1 \rceil + n \cdot \theta_1 - n \cdot \theta_1 }{n+1} .
% $$

% \begin{wrapfigure}{r}{0.25\textwidth}
%     \centering
%     \includegraphics[width=0.5\textwidth]{img/support_distribution.pdf}
%     \caption{The support threshold $\theta_1$ in relation to the stream threshold $s$ for $n=4$.}
%     \label{fig:thresholds}
% \end{wrapfigure}



\subsection*{(c): Updating both, the Stream Database
and the Support Threshold}

Finally, we consider the most general case of a
change, where streams are added to the database and
the support threshold is increased. Again, we note
that the change in the stream thresholds that is
induced by the change determines whether the sets of
matching queries are in a containment relation, i.e.,
whether the change leads only to frequent queries
being no longer frequent, but not to additional
frequent queries. We capture this property as follows:
\begin{lemma}
    \label{lemma:update}
    Let $D_0=\{S_1,\dots,S_d\}$ and
    $D_1=D_0 \cup \{S_{d+1}, \dots, S_{d+k}\}$
    be stream databases, and $\theta_0$ and $\theta_1$
    with $\theta_0 \leq \theta_1$ be support
    thresholds for $D_0$ and $D_1$, respectively.
    Then, it holds that $s(|D_1|,\theta_1) \geq
    s(|D_0|,\theta_0) + k \implies Q_{D_1}^{\theta_1}
    \subseteq Q_{D_0}^{\theta_0}$.
\end{lemma}
\textit{Proof.}
Let $q_1$ be a query in $Q_{D_1}^{\theta_1}$. Recall
that $M_+(q_1,D_1)$ is the set of streams in $D_1$
that match $q_1$, $M_+(q_1,D_0)$ is the set of streams
in $D_0$ that match $q_1$, and $M_+(q_1,D_1\backslash
D_0)$ is the set of newly added streams that match
$q_1$. Let $|M_+(q_1,D_1\backslash D_0)|= j$ be the
number of these streams, for which it holds that $j\in
[0,k]$. Then, it suffices to show that $|M_+(q_1,D_0)|
> s(|D_0|,\theta_0)$. As before, we define $s_0 =
s(|D_0|,\theta_0)$ and $s_1 = s(|D_1|,\theta_1)$.
Then, we derive the following bounds for the
cardinality of the aforementioned sets:
$$|M_+(q_1,D_0)| = |M_+(q_1,D_1)| -
|M_+(q_1,D_1\backslash D_0)| \geq s_1 - j \geq s_0 + k
- j \geq s_0.$$
% Then $|M_2_+| \geq s_2 > s_1$.
% If $\{S_{n+1}, \dots, S_{n+k}\}\subseteq M_+(q_1,D_1)$.
% $$|M_+(q_1,D_1)\backslash \{S_{n+1}, \dots, S_{n+k}\}| = s_1 - k
% \geq s_0 + k - k = s_0.$$
% $$|M_2_+\backslash \{S\}| \geq \lceil (n+1) \cdot \theta_1 \rceil - 1
% \geq \lceil n \cdot \theta_1 \rceil = s.$$
% Then $M_+(q_1,D_1)\backslash \{S_{d+1}, \dots, S_{d+k}\} = M_+(q_1,D_0)$ and supp$(q_1,D_0) \geq \frac{|M_+(q_1,D_0)|}{|D_0|} \geq \theta_0$.
% If  $\exists S \in \{S_{d+1}, \dots, S_{d+k}\}$ s.t.
%$S\not\in M_+(q_1,D_1)$:
% If $\{S_{n+1}, \dots, S_{n+k}\}\not\subseteq M_+(q_1,D_1)$.
% Then $M_+(q_1,D_1) \subseteq D_0$.
% $$\frac{|M_+(q_1,D_1)|}{|D_0|} \geq \frac{\lceil (n+k) \cdot \theta_1 \rceil}{n}
% \geq \frac{\lceil n \cdot \theta_0 \rceil}{n} \geq \frac{n\cdot\theta_0}{n} = \theta_0.$$
% Then supp$(q_1,D_0) \geq \theta_0$ and $q_1 \in Q_0$.
Hence, indeed, it holds that $Q_{D_1}^{\theta_1}
\subseteq Q_{D_0}^{\theta_0}$. \hfill  \qed

We note that \autoref{lemma:update} provides a
generalization of \autoref{prop:db_change_inclusion}.
That is, if $\theta_1 = \theta_0$, i.e., streams are
added to the database, but the support threshold
remains unchanged, due to the definition of the stream threshold, $s_0+k$ is an
upper bound for $s_1$. As such, the condition of the implication on the
containment relation becomes $s_1 = s_0+k$.

The above lemma gives a condition, i.e., $s_1\geq
s_0+k$, to conclude on a containment relation for the
sets of matching queries under an update of the stream
database, regardless of the streams that are added.

Considering the general problem context as visualized
in \autoref{fig:matching_sets}, one might assume that
$Q_{D_1}^{\theta_1}\subseteq Q_{D_0}^{\theta_0}
\implies Q_{D_1}^{\theta_0}\subseteq
Q_{D_0}^{\theta_0}$ or $Q_{D_1}^{\theta_1}\subseteq
Q_{D_0}^{\theta_1}$. That is, if the set of matching
queries is only reduced by changes of potentially
both, the stream database and the support threshold,
the containment relation is expected to hold also if
only one of the changes is considered.
However, this is actually not the case, which we
illustrate with a counter example in
\autoref{tab:example}.

\begin{table}[h!]
  \centering
  \begin{tabular}{@{}cccccccc@{}}
  \toprule
  $d$ & $d+1$ & $\theta_0$ & $\theta_1$ &
  $s(d,\theta_0)$ & $s(d+1,\theta_0)$ &
  $s(d,\theta_1)$ & $s(d+1,\theta_1)$ \\ \midrule
  99  & 100   & 0.8        & 0.9        & 80              & 80                & 90              & 90                \\ \bottomrule
  \end{tabular}%
  \caption{Stream thresholds under an example change
  in the stream database and the support threshold.}
  \label{tab:example}
  \end{table}

Here, we consider the setting of a stream database of
size $d=99$, for which one stream is added to obtain a
database of size $d+1=100$. Moreover, we consider a
change of the support threshold from $\theta_0=0.8$ to
$\theta_1=0.9$. For this setting, it holds that
$s(d+1,\theta_0) = s(d,\theta_0)$ as well as
$s(d+1,\theta_1) = s(d,\theta_1)$, i.e., the stream threshold does not change, if
though a stream was added to the database. Hence, without information on this
stream, we cannot
conclude that $Q_{D_1}^{\theta_0} \subseteq
Q_{D_0}^{\theta_0}$ or
$Q_{D_1}^{\theta_1} \subseteq Q_{D_0}^{\theta_1}$.
However, the change of both, the stream database and the support threshold,
actually leads to a change in the stream threshold, i.e., $s(d+1,\theta_1) \geq
s(d,\theta_0)$. Based thereon, using \autoref{lemma:update}, we can conclude on
the containment relation for the sets of matching queries, $Q_{D_1}^{\theta_1}
\subseteq Q_{D_0}^{\theta_0}$.

% We can now take a look at the probability of the different cases for the stream thresholds $s_2$ and $s_1$,
% namely $s_2 = s_1$ and $s_2 \geq s_1+1$.

% With $s_1 = \lceil n \cdot \theta_1 \rceil$ and $s_2 = \lceil (n+1) \cdot \theta_2 \rceil$ we know the boundaries for
% $\theta_1$ and $\theta_2$ in relation to the stream thresholds $s_1$ and $s_2$:

% \begin{align*}
%      \frac{s_1-1}{n} &< \theta_1 \leq \frac{s_1}{n} \\
%     \frac{s_2}{n+1} < &\theta_2 \leq \frac{s_2}{n+1}.
% \end{align*}

% Subtracting the two inequalities we obtain:
% $$\frac{s_2-1}{n+1} - \frac{s_1}{n} < \theta_2 - \theta_1 < \frac{s_2}{n+1} - \frac{s_1-1}{n}.$$

% Now we can look at the two cases:\\
% $s_2 = s_1$:
% \begin{align}
%      \frac{s_1-1}{n} - \frac{s_1}{n} &< \theta_2 - \theta_1 < \frac{s_1}{n+1} - \frac{s_1}{n+1} \nonumber\\
%      \implies -\frac{1}{n+1} - \frac{s_1}{n(n+1)} &< \theta_2 - \theta_1 < \frac{1}{n} - \frac{s_1}{n(n+1)} \nonumber \\
%       \implies 0 &< \theta_2 - \theta_1 < \frac{1}{n} - \frac{s_1}{n(n+1)}. \label{eq:case1}
%     \end{align}

%   $s_2 \geq s_1+1: s_2=s_1+k+1, k\in[0,1,\dots,n-1]$:
%   \begin{align}
%     \frac{s_1+k}{n+1} - \frac{s_1}{n} &< \theta_2 - \theta_1 < \frac{s_2}{n+1} - \frac{s_1-1}{n} \nonumber \\
%     \implies \frac{k}{n+1} - \frac{s_1}{n(n+1)} &< \theta_2 - \theta_1 < \frac{1}{n} + \frac{k+1}{n+1} - \frac{s_1}{n(n+1)} \label{eq:case2} \\
%   k=0:  \implies 0 &< \theta_2 - \theta_1 < \frac{1}{n} + \frac{1}{n+1} - \frac{s_1}{n(n+1)}. \label{eq:case2k0}
%   \end{align}

%   To obtain all possible combinations of $\theta_1$ and $\theta_2$ we have to consider all possible values for $s_1$ and $k$.
%   $$
%   \sum_{s_1=1}^{n} \left( \underbrace{\frac{1}{n} - \frac{s_1}{n(n+1)}}_{\autoref{eq:case1}} +
%    \underbrace{\frac{1}{n} + \frac{1}{n+1} - \frac{s_1}{n(n+1)}}_{\autoref{eq:case2k0}} +
%   \sum_{k=1}^{n+1-s_1} \underbrace{\frac{1}{n} - \frac{s_1}{n(n+1)}}_{\autoref{eq:case2}} \right) = \frac{1}{2} + n - \frac{1}{2} = n.
%   $$

%   For the probabilities of the different cases we obtain:
%   \begin{align*}
%     &\mathbb{P}[s_2 = s_1] = \frac{1}{2n} \\
%     &\mathbb{P}[s_2 \geq s_1+1] = \frac{2n-1}{2n}.
%   \end{align*}

